% !TEX root = ../main.tex
\lecture{10}{2026-11-16}{Taylorpolynome}

\subsubsection{Konvexitätserhaltung}
\begin{definition}
    [Konvexitätserhaltend] sind f und g auf \textit{I}
    \begin{itemize}
        \item (streng) konvex, sind f + g und max\{f,g\} auf I
        (streng) konvex
        \item (streng) konkav, ist f + g und min\{f,g\} auf I
        (streng) konkav
        \item max einer konvexen/konkaven und einer streng konvexen/konkaven
        Funktion ist auch konvex/konkav aber nicht streng
        konvex
    \end{itemize} 
\end{definition} 


\subsubsection{Wendepunkt}
\begin{definition}
    [Wende und Sattelpunkt] ist eine Funktion auf einem
    Interval konvex und auf dem folgenden konkav muss es am
    Schnittpunkt dieser Intervale einen Wendepunkt oder
    Sattelpunkt befinden. Sattelpunkt wenn f differenzierbar: $f'(x_0) = 0$
\end{definition}

\begin{theorem}
    [Kriterium Wendepunkt] wenn f einen Wendepunkt hat und zweimal differenzierbar in
    $x_0$ gilt:$f''(x_0) = 0 $ aber wenn f''(x) = 0 gilt
    \textbf{nicht} x als Wendepunkt (Bsp. $x^4$)
\end{theorem}

\subsubsection{Symmetrie und Periodizität}

\begin{definition}
    [Symmetrie] Gilt $f: D \rightarrow Z und alle x \in D$
    \begin{itemize}
        \item f(-x) = f(x) spiegelsymmetrisch um y-achse
        \item f(-x) = -f(x) punktsymmetrisch um Ursprung
        \item f(x) = f(x + T) mit T > 0 ist periodisch mit
        Periode T
    \end{itemize} 

\end{definition}

\subsection{Taylor-Polynom}
Besteht aus dem Versuch eine Funktion über ihre Werte zu
approximieren, in dem man mit einer Tangente an einem bekannten
Punkt anfängt (Wert an der Ursprungsfunktionstelle ist gleich
Tangentenstelle und der ersten Ableitung). Danach ist eine
Funktion gesucht die weiterhin diese Eigenschaften hat, aber
auch die 2. Ableitung bis n-te Ableitung entspricht. So
Annäherungsweise wird die Ursprungsfunktion über Polynome
hergestellt mit einem minimierenden Fehler (r). Also wenn
Funktionsvorschrift unbekannt aber Eigenschaften bekannt kann
diese über Taylorpolynome 

\begin{definition}
    [Taylor-Polynom mit Restglied] Sei $f$ $(n+1)$-mal
    differenzierbar auf einem Intervall, das $x_0$ und $x$ enthält.
    Das Taylor-Polynom von $f$ zum Entwicklungspunkt $x_0$ und zur
    Ordnung $n$ ist
    \[
        T_n(x) = \sum_{k=0}^{n} \frac{f^{(k)}(x_0)}{k!}(x-x_0)^k.
    \]
    Nach dem Satz von Taylor gilt
    \[
        f(x) = T_n(x) + R_n(x), \qquad
        R_n(x) = \frac{f^{(n+1)}(\xi)}{(n+1)!}(x-x_0)^{n+1},
    \]
    wobei $\xi$ zwischen $x_0$ und $x$ liegt. $R_n$ heisst
    Restglied in der Lagrange-Form.
\end{definition}
 